\documentclass[10pt,twocolumn]{article}

\usepackage[margin=0.75in,columnsep=0.28in]{geometry}
\usepackage{amsmath,amssymb}
\usepackage{zref-savepos}

% Record the PDF position and TeX dimensions of every marked expression.
\newsavebox{\mathbox}
\newwrite\mathcoords
\immediate\openout\mathcoords=mathcoords.csv
\newcounter{mathid}
\immediate\write\mathcoords{id,page,x,y,width,height,depth}

\newcommand{\recordmath}[1]{%
    \stepcounter{mathid}%
    \sbox{\mathbox}{$#1$}%
    \edef\mathlabel{math-\themathid}%
    % This is important: without it, a marker at the start of a paragraph is
    % recorded in vertical mode, before indentation and line placement.
    \leavevmode
    \zsavepos{\mathlabel}%
    \immediate\write\mathcoords{%
        \themathid,
        \thepage,
        \zposx{\mathlabel},
        \zposy{\mathlabel},
        \number\wd\mathbox,
        \number\ht\mathbox,
        \number\dp\mathbox
    }%
    \usebox{\mathbox}%
}

% Use display style while retaining the same recording mechanism.
\newcommand{\recorddisplay}[1]{\recordmath{\displaystyle #1}}

\AtEndDocument{\immediate\closeout\mathcoords}

\title{A Synthetic Research Paper for Testing Math Bounding Boxes}
\author{MathBox Test Suite}
\date{}

\begin{document}

\maketitle

\begin{abstract}
This synthetic paper contains representative mathematical constructions at
different locations in a two-column document. It is intended for testing the
detection and placement of equation bounding boxes rather than presenting a
scientific result.
\end{abstract}

\section{Inline and display mathematics}

An inline equation such as \recordmath{E=mc^2} should be detected without
including the surrounding sentence. Multiple expressions can also occur on
one line: \recordmath{a=b}, \recordmath{x_1+x_2=1}, and
\recordmath{f(t)=e^{-t}} are three separate inline expressions.

An unnumbered display equation tests centered mathematics:
\[
    \recorddisplay{f(x)=x^2+2x+1.}
\]

The following numbered equation contains subscripts and superscripts:
\begin{equation}
    \recorddisplay{y_i=\beta_0+\beta_1x_i+\varepsilon_i.}
    \label{eq:regression}
\end{equation}

Equation~\eqref{eq:regression} is followed by ordinary prose so that the
vertical extent of the equation can be distinguished from adjacent text.

\section{Fractions and integrals}

Fractions of different heights are useful for checking both the height above
the baseline and the depth below it:
\begin{equation}
    \recorddisplay{p(x)=\frac{x^2+1}{x^3-x+1}.}
    \label{eq:fraction}
\end{equation}

Here is an integral with limits and a nested fraction:
\begin{equation}
    \recorddisplay{I_n=\int_{0}^{1}\frac{x^n}{1+x^2}\,\mathrm{d}x.}
    \label{eq:integral}
\end{equation}

Inline fractions provide a different baseline case: the value
\recordmath{r_i=\frac{u_i-v_i}{1+v_i^2}} appears in the middle of this sentence.

\section{Aligned equations}

The \texttt{align} environment produces several numbered rows whose alignment
points do not necessarily share the same horizontal extent:
\begin{align}
    \recorddisplay{\nabla\!\cdot\mathbf{E}} &={}
        \recorddisplay{\frac{\rho}{\varepsilon_0}},
        \label{eq:gauss-electric}\\
    \recorddisplay{\nabla\!\cdot\mathbf{B}} &={}
        \recorddisplay{0},
        \label{eq:gauss-magnetic}\\
    \recorddisplay{\nabla\!\times\mathbf{E}} &={}
        \recorddisplay{-\frac{\partial\mathbf{B}}{\partial t}}.
        \label{eq:faraday}
\end{align}

The next display is deliberately placed near the bottom of the first page.
The flexible space pushes it down while keeping the source easy to modify.

\vfill
\begin{equation}
    \recorddisplay{q_{\mathrm{bottom},1}=\sum_{k=1}^{n}\frac{k}{n^2}.}
    \label{eq:bottom-one}
\end{equation}

\clearpage

\section{Matrices and systems}

A matrix tests tall delimiters and entries extending above and below a normal
text baseline:
\begin{equation}
    \recorddisplay{\mathbf{A}=
    \begin{pmatrix}
        a_{11} & a_{12} & a_{13}\\
        a_{21} & a_{22} & a_{23}\\
        a_{31} & a_{32} & a_{33}
    \end{pmatrix}.}
    \label{eq:matrix}
\end{equation}

A piecewise definition is another useful tall construction:
\begin{equation}
    \recorddisplay{g(x)=
    \begin{cases}
        x^2, & x\geq 0,\\
        -x,  & x<0.
    \end{cases}}
    \label{eq:piecewise}
\end{equation}

The following aligned system is unnumbered:
\begin{align*}
    \recorddisplay{2x+3y-z} &={} \recorddisplay{7},\\
    \recorddisplay{x-y+4z} &={} \recorddisplay{2},\\
    \recorddisplay{3x+2y+2z} &={} \recorddisplay{9}.
\end{align*}

\section{Long and nested expressions}

This equation occupies much of a column's width:
\begin{equation}
    \recorddisplay{\mathcal{L}(\theta)
    =-\sum_{i=1}^{N}\left[y_i\log p_i
      +(1-y_i)\log(1-p_i)\right].}
    \label{eq:likelihood}
\end{equation}

Nested constructions test whether a box includes every level:
\begin{equation}
    \recorddisplay{h(x)=\sqrt{1+\frac{x_1^2}{1+\frac{x_2^2}{1+x_3^2}}}.}
    \label{eq:nested}
\end{equation}

Two displays can also appear close together:
\[
    \recorddisplay{u(t)=u_0+at,}
\]
\[
    \recorddisplay{s(t)=s_0+u_0t+\frac{1}{2}at^2.}
\]

The intervening prose helps reveal whether nearby boxes overlap or are merged.
For an inline stress test, consider \recordmath{\alpha_i},
\recordmath{\beta_{ij}}, \recordmath{\gamma^{(k)}}, and
\recordmath{T_{\mu\nu}} within the same line of text.

\vfill
\begin{equation}
    \recorddisplay{q_{\mathrm{bottom},2}=\lim_{m\to\infty}
    \left(1+\frac{1}{m}\right)^m=e.}
    \label{eq:bottom-two}
\end{equation}

\clearpage

\section{Page-boundary stress tests}

The final page combines several forms that are common in technical papers.
First, a Fourier transform contains an infinite integral and complex
exponential:
\begin{equation}
    \recorddisplay{\widehat{f}(\omega)=\int_{-\infty}^{\infty}
    f(t)e^{-i\omega t}\,\mathrm{d}t.}
    \label{eq:fourier}
\end{equation}

Next, an optimization problem uses constraints and multiple aligned rows:
\begin{align}
    \recorddisplay{\underset{\mathbf{x}}{\operatorname{minimize}}}\quad
        & \recorddisplay{\frac{1}{2}\mathbf{x}^{\mathsf T}\mathbf{Q}\mathbf{x}
          +\mathbf{c}^{\mathsf T}\mathbf{x}} \\
    \text{subject to}\quad
        & \recorddisplay{\mathbf{A}\mathbf{x}=\mathbf{b}},
        \label{eq:optimization}\\
        & \recorddisplay{x_i\geq 0 \quad (i=1,\ldots,n)}.
\end{align}

A short expression \recordmath{P(A\mid B)=P(A\cap B)/P(B)} sits between
displays. The next equation tests large operators and deeply placed subscripts:
\begin{equation}
    \recorddisplay{Z_N=\prod_{j=1}^{N}\left(
        \sum_{k_j=0}^{\infty}e^{-\lambda_j k_j}
    \right).}
    \label{eq:operators}
\end{equation}

The last equation is pushed toward the bottom of the final page. This is useful
for detecting clipping and off-by-one-page errors in coordinate conversion.

\vfill
\begin{equation}
    \recorddisplay{\boxed{R_{\mathrm{final}}
      =\int_{0}^{2\pi}\!\int_{0}^{1}r^2\sin\theta
      \,\mathrm{d}r\,\mathrm{d}\theta}}
    \label{eq:final}
\end{equation}

\end{document}
